% Reader question: When a process crashes mid-task, what breaks and what
%   survives across the three layers a transaction touches?
% Claim: structural authority breaks at the crash and is re-issued to a
%   successor, but the evidence chain and the pre-funded escrow both span the
%   crash unbroken.
% Evidence: the worked example minute-by-minute timeline (capability, claim,
%   crash at minute 12, successor capability at minute 13, settlement at
%   minute 24) from this chapter worked example [internal].
% Must distinguish: which layer the crash actually interrupts (L1 structural
%   authority only) from the two layers that bridge it unbroken (L2 evidence,
%   L3 escrow).
% Grammar chosen: three-lane wall-clock timeline, one shared minute axis.
%   Rejected: tinted background strata for each lane, because a tinted band
%   reads as a sticker on a white page and the lane identity is already
%   carried by its own label and rule.
% Five-second test: a reader names the one layer that breaks at the crash and
%   the two that do not.
\begin{figure}[H]
\centering
\pdfigurehue{pdgold}%
\begin{tikzpicture}[pd figure,x=1cm,y=1cm]
  \def\xSt{1.60}  \def\xEnd{8.80}
  \def\mB{2.30}  \def\mC{3.40}  \def\mCrash{5.60}  \def\mSucc{6.60}
  \def\yL{3.30}  \def\yE{1.90}  \def\yS{0.50}  \def\yax{-1.10}
  \node[pd title,anchor=west,text width=8.4cm] at (\xSt,4.30)
    {a crash breaks the process, not the task evidence or escrow};
  \node[pd label,anchor=east] at ({\xSt-.14},\yL) {L1 scope};
  \node[pd label,anchor=east] at ({\xSt-.14},\yE) {L2 evidence};
  \node[pd label,anchor=east] at ({\xSt-.14},\yS) {L3 escrow};
  \foreach \y in {\yL,\yE,\yS} {\draw[pd hairline] (\xSt,\y)--(\xEnd,\y);}
  % ---- L1: structural authority ends at the crash, re-issued after --------
  \draw[pd rule,line width=1.4pt] (\xSt,\yL)--(\mCrash,\yL);
  \draw[pd focus rule,line width=1.4pt] (\mSucc,\yL)--(\xEnd,\yL);
  \node[pd datum] at (\xSt,\yL) {};
  \node[pd datum] at (\mC,\yL) {};
  \node[pd breach datum] at (\mCrash,\yL) {};
  \node[pd focus datum] at (\mSucc,\yL) {};
  \draw[pd breach rule,line width=1pt] ({\mCrash-.11},{\yL-.20})--({\mCrash+.11},{\yL+.20});
  \draw[pd breach rule,line width=1pt] ({\mCrash+.02},{\yL-.20})--({\mCrash+.24},{\yL+.20});
  \node[pd label,anchor=south] at (\xSt,{\yL+.20}) {capability};
  \node[pd label,anchor=south] at (\mC,{\yL+.20}) {claim};
  \node[pd breach label,anchor=south east] at ({\mCrash-.10},{\yL+.20}) {crash};
  \node[pd ready label,anchor=south west] at ({\mSucc+.10},{\yL+.20}) {successor cap};
  % ---- L2: the evidence chain bridges the crash unbroken -------------------
  \draw[pd focus rule,line width=1.15pt] (\mC,\yE)--(\mCrash,\yE);
  \draw[pd focus rule,line width=1.15pt] (\mSucc,\yE)--(\xEnd,\yE);
  \foreach \i in {0,...,4} {\node[pd focus datum,inner sep=1.4pt] at ({\mC+\i*.675},\yE) {};}
  \foreach \i in {0,...,2} {\node[pd focus datum,inner sep=1.4pt] at ({\mSucc+\i*.82},\yE) {};}
  \draw[pd focus arrow,dashed] (\mCrash,{\yE+.10}) .. controls ({\mCrash+.20},{\yE+.55}) and ({\mSucc-.20},{\yE+.55}) .. (\mSucc,{\yE+.10});
  \node[pd ready label,anchor=south] at ({(\mCrash+\mSucc)/2},{\yE+.58}) {salvage parent};
  \node[pd note,anchor=north] at (\mC,{\yE-.20}) {$h_1\ldots h_5$};
  \node[pd note,anchor=north] at (\mSucc,{\yE-.20}) {$h_6\ldots h_8$};
  % ---- L3: pre-funded escrow spans the full task lifetime -------------------
  \draw[pd breach rule,line width=2.4pt] (\mB,\yS)--(\xEnd,\yS);
  \node[pd breach datum] at (\mB,\yS) {};
  \node[pd focus datum] at (\xEnd,\yS) {};
  \node[pd breach label,anchor=south] at (\mB,{\yS+.20}) {bond \$6};
  \node[pd ready label,anchor=south] at (\xEnd,{\yS+.20}) {settle to $P$};
  % ---- minute axis --------------------------------------------------------
  \draw[pd rule,-{Stealth[length=1.8mm,width=1.1mm]}] (\xSt,\yax)--({\xEnd+.15},\yax);
  \node[pd label,anchor=west] at ({\xEnd+.20},\yax) {minute};
  \foreach \x/\lab in {\xSt/0,\mB/1,\mC/3,\mCrash/12,\mSucc/13,\xEnd/24} {
    \draw[pd tick] (\x,{\yax-.07})--(\x,{\yax+.07});
    \node[pd label,anchor=north] at (\x,{\yax-.12}) {\lab};}
\end{tikzpicture}
\caption{The worked example is a continuity braid over one task. Structural
authority belongs to a process and visibly breaks at minute 12; the
successor receives a new capability at minute 13. The evidence chain does
not break: its salvage-parent edge binds $h_1\ldots h_5$ to $h_6\ldots h_8$.
The escrow rail spans the entire task and settles at minute 24. Identity is
disposable, while attribution and pre-funded cleanup exposure persist.
[internal]}
\label{fig:worked-example}
\end{figure}
